\documentclass[xcolor={dvipsnames,table}]{beamer}

%  Include preamble
%
% All Thesis package dependencies
%

%% --- Core Engine & Language ---
\usepackage[british]{babel}

% REMOVED: inputenc and fontenc. 
% fontspec handles encoding natively in LuaLaTeX/XeLaTeX.

% --- Fonts & Unicode Support ---
\usepackage{fontspec}
\setmainfont{TeX Gyre Termes}  % Professional Serif (Times-like)
\setsansfont{TeX Gyre Heros}   % Professional Sans (Helvetica-like)

% --- Typography ---
\usepackage{microtype}
% Disable specific microtype features that clash with LuaLaTeX/fontspec
\microtypesetup{expansion=false, kerning=false}

% --- Coding & Formalism ---
\usepackage{amsmath}
\usepackage{listings}
\usepackage{amssymb}
\usepackage{pifont}
% \usepackage{dingbat}

% --- Proof trees ---
\usepackage{ebproof}
% Custom style to make the horizontal lines look better
% \ebproofset{line thickness=0.6pt, separation=0.5em}

% --- Referencing & Utilities ---
\hypersetup{hidelinks}
\usepackage[
    style=authoryear, 
    backend=biber,    % Recommended backend over bibtex
    natbib=true,      % Supports \citet and~\citep
]{biblatex}

% 1. Replace the error-causing line with these:
\DeclareFieldFormat{citebackslash}{#1}
\renewcommand{\bibopenparen}{\addthinspace[}
\renewcommand{\bibcloseparen}{]}

% 2. Ensure the delimiter between name and year is correct:
\renewcommand*{\nameyeardelim}{\addcomma\space}

\addbibresource{sources.bib}

\usepackage[capitalise]{cleveref}
\usepackage{xspace}
\usepackage{todonotes}
\usepackage{lipsum}
\usepackage{booktabs}

% --- Margin & Formatting Tools ---
\usepackage{environ}
\usepackage{ragged2e}
\usepackage{changepage}

\usepackage[most]{tcolorbox}
\usepackage{tabularx}
\usepackage{tikz}
% Import positioning and arrows libraries for our custom overlays
\usetikzlibrary{positioning, arrows.meta, tikzmark}
%
% All Thesis custom macros
% and redefinitions
%

\renewcommand{\bfdefault}{bx}
\DeclareTextFontCommand\textsb{\libertineSB}

\Crefname{section}{\S}{\S\S}

% TikZ helper to create "invisible" anchors in normal text or code blocks
\newcommand{\tikztarget}[2]{\tikz[remember picture,baseline=(#1.base)] \node[inner sep=0pt] (#1) {#2};}

\newcommand\myhrulefill[1]{\leavevmode\leaders\hrule height#1\hfill\kern0pt}

\newcommand{\RED}[1]{\textcolor{red}{#1}}
\newcommand{\ie}{\emph{i.e.,}}
\newcommand{\eg}{\emph{e.g.,}}
\newcommand{\etal}{\emph{et~al.}}


\newcommand{\Yes}{\ding{51}}
\newcommand{\No}{\ding{55}}

\renewcommand{\bfdefault}{b}

\newcommand{\While}{\textbf{While}\xspace}
\newcommand{\Num}{\textbf{Num}\xspace}
\newcommand{\Var}{\textbf{Var}\xspace}
\newcommand{\Aexp}{\textbf{Aexp}\xspace}
\newcommand{\Bexp}{\textbf{Bexp}\xspace}
\newcommand{\Stm}{\textbf{Stm}\xspace}
\newcommand{\State}{\textbf{State}\xspace}
\newcommand{\Truth}{\textbf{T}\xspace}

\newcommand{\nb}[3]{
    {\colorbox{#2}{\bfseries\sffamily\scriptsize\textcolor{white}{#1}}}{
        \textcolor{#2}{\sf\small\textit{#3}}}}
        
% \newcommand{\hampus}[1]{\nb{Hampus}{magenta}{#1}}
% \newcommand{\elias}[1]{\nb{Elias}{Orchid}{#1}}
% \newcommand{\andrea}[1]{\nb{Andrea}{Brown}{#1}}
% \newcommand{\missingcitation}{\nb{Citation missing}{ForestGreen}{}\xspace}
% \newcommand{\question}[1]{\nb{Question}{RoyalBlue}{#1}\xspace}
% \newcommand{\warning}[1]{\nb{Warning}{Red}{#1}\xspace}
% \newcommand{\TODO}[1]{\nb{TODO}{Orange}{#1}\xspace}
% \newcommand{\note}[1]{\nb{Note}{Black!60}{#1}\xspace}
% \newcommand{\suggestion}[1]{\nb{Suggestion}{BlueGreen}{#1}\xspace}

\newcommand{\elias}[1]{}
\newcommand{\hampus}[1]{}
\newcommand{\missingcitation}{}
\newcommand{\question}[1]{}
\newcommand{\warning}[1]{}
\newcommand{\TODO}[1]{}
\newcommand{\suggestion}[1]{}
\renewcommand{\todo}[1]{}


% --- lean-style.tex ---

% Define the lean4 style
\lstdefinestyle{lean4}{
  language=bash, % Fallback
  basicstyle=\normalsize\ttfamily\color{black},
  keywordstyle=\ttfamily\bfseries\color{blue!70!black},
  commentstyle=\itshape\color{green!40!black},
  showstringspaces=false,
  stringstyle=\color{purple},
  % Clear all default strings and only allow double quotes
  morestring=[b]", 
  % Explicitly tell listings to ignore single quotes as delimiters
  deletestring=[d]',
  deletestring=[b]',
  % Ensure the tick/prime is treated as a normal character
  alsoletter={'},
  numberstyle=\footnotesize\sffamily\color{gray},
  tabsize=2,
  % xleftmargin=1ex,
  showstringspaces=false,
  % Lean 4 Keywords
  morekeywords={
    def, theorem, lemma, example, open, import, 
    variable, variables, instance, class, structure, 
    inductive, syntax, macro, deriving, if, then, else, 
    let, have, show, by, match, with, infix, notation, namespace, end,
    apply, intro, simp, grind, exact, split_ifs, repeat, rw,  rfl, subst,
    contradiction, where, Prop, constructor, exists, specialize, linarity,
    cases, induction, generalizing, rename_i, linarith, prefix, infix,
    infixl, infixr, trivial, obtain, using, and_intros
  },
% Comments
  morecomment=[l]{--},
  morecomment=[s]{/-}{-/},
  % Symbols and Delimiters
  escapeinside=§§,
  % Unicode Handling for listings
  % Even with LuaLaTeX, 'literate' ensures symbols look perfect in code
  extendedchars=true,
  literate=
    {→}{{$\ →\ $}}3
    {↦}{{$\ \mapsto\ $}}3
    {⇒}{{$⇒\ $}}2
    {∀}{{$\forall\ $}}2
    {∃}{{$\exists\ $}}2
    {λ}{{$\lambda\ $}}2
    {:=}{{$:=\ $}}2
    {::=}{{$::=\ $}}3
    {∘}{{$\circ\ $}}2
    {|}{{$\mid\ $}}2
}

% Choose a theme (optional, but makes it look pretty)
% \usetheme{Madrid}
\usetheme{metropolis}
\usecolortheme{default}

\title{Mechanizing Proofs from a Course on Programming Language Semantics in Lean}
\subtitle{Bachelors Thesis Presentation}
\author{Hampus Toft}
\date{\today}

\begin{document}

% Slide 1: Title Page
\begin{frame}
  \titlepage
\end{frame}

% Slide 2: Table of Contents
\begin{frame}{Agenda}
  \begin{itemize}
    \item What is this thesis about?
    \item How do you even compare Lean and pen-and-paper?
    \item Examples of the process
    \item Results of the mechanization process
    \item Conclusion
  \end{itemize}
\end{frame}

% Slide 3: Introduction
\begin{frame}
  \section*{Introduction}
  What is the meaning of: \\
  \emph{"Mechanizing Proofs from a Course on Programming Language Semantics in Lean"}
  
  \begin{block}<2->{Mechanization}
    Formal proof in a computer proof assistant.
  \end{block}

  \begin{block}<3->{Programming Language Semantics}
    The meaning behind the code.
  \end{block}

\end{frame}

% Slide 4: Introduction (Nielson)
\begin{frame}
  \section*{Nielson \& Nielson's Framework}
  Case study using "\textbf{Semantics with Applications: An Appetizer}". \\
  With the focus on:
  \begin{itemize}
    \item<1-> Chapter 1: The \While language
    \item<2-> Chapter 2: Operational Semantics
      \begin{itemize}
        \item<3-> Natural Operational Semantics
        \item<4-> Structural Operational Semantics
      \end{itemize}
  \end{itemize}
\end{frame}

% Slide 5: Introduction (Lean)
\begin{frame}
  \section*{The Lean 4 Proof Assistant}
  Uses "Dependent Type Theory". \\
  With features such as:
  \begin{itemize}
    \item<1-> Inductive Types
    \item<2-> Tactics
    \item<3-> Strict Correctness
  \end{itemize}
\end{frame}

% Slide 6: Research Questions
\begin{frame}
  \section*{Research Questions}
  \begin{itemize}
    \item<1-> \textbf{RQ1:} Do the proof techniques presented in "\emph{Semantics with Applications: An Appetizer}"
    map directly to Lean's internal mechanisms, or do they require significant
    structural reformulation?

    \item<2-> \textbf{RQ2:} What implicit assumptions or "pedagogical shortcuts"
    present in "\emph{Semantics with Applications: An Appetizer}" must be made explicit to satisfy the Lean proof assistant?
  \end{itemize}
\end{frame}

% Slide 7: Method
\begin{frame}
    \section{Method}
    \begin{enumerate}
        \item Find a something to mechanize
        \item Then mechanize it
        \item Put side-by-side
        \item Repeat Steps 1-3
        \item \dots
        \item Profit?
    \end{enumerate}
\end{frame}

% Slide 8: Mechanization of Aexp
\begin{frame}[fragile]{The 6-Step Program to Profit}
    \begin{columns}[T]
        \begin{column}{0.55\textwidth}
            \begin{block}{Method}
\begin{enumerate}
    \item<1-> Find something to mechanize
    \item<2-> Then mechanize it
    \item<3-> Put them side-by-side
    \item<5-> Repeat Steps 1-3
\end{enumerate}
            \end{block}
        \end{column}
        \begin{column}{0.45\textwidth}
            \begin{block}{Checklist}
\begin{itemize}
    \item Inductive types \No
    \item Semantic Functions \No
    \item Inference Rules \No
    \item Proofs and Lemmas \No
\end{itemize}
            \end{block}
        \end{column}
    \end{columns}

    % \vspace{0.5cm}
    
    % We wrap the side-by-side view inside an overprint or conditional block 
    % so it only shows up from Slide 1 onwards, changing state dynamically.
    \begin{columns}[T]
        
        % --- LEFT COLUMN: Lean Code ---
        \begin{column}{0.50\textwidth}
            % Show right column on slide 2 (isolated) and slides 3+ (side-by-side)
            \begin{block}<2->{Lean 4 Formalization}
                \begin{lstlisting}[style=lean4]
inductive Aex§\tikzmarknode{ex1_ln1}{p}§ where
  | num (n : Num) §\tikzmarknode{ex1_ln2}{}§
  | var (x : Var)
  | add (a1 a2 : Aexp)
  | mul (a1 a2 : Aexp)
  | sub (a1 a2 : Aexp)
                \end{lstlisting}
            \end{block}
        \end{column}

        % --- RIGHT COLUMN: Original Text Target---
        \begin{column}{0.48\textwidth}
            % Only show left column from slide 1, 3, and 4 (Skip slide 2 to isolate Lean)
            \only<1,3->{
                \begin{block}{Original Definition}
\[
\begin{array}{lcl}
    \\[0.3em]
    \tikztarget{ex1_txt1}{a} &::=& \tikztarget{ex1_txt2}{n} \\ 
        &\mid& x \\ 
        &\mid& a_1 + a_2 \\ 
        &\mid& a_1 \star a_2 \\ 
        &\mid& a_1 - a_2 \\
\end{array}
\]
                \end{block}
            }
        \end{column}
        
    \end{columns}

    % --- TIKZ OVERLAY LAYER ---
    % This executes "on top" of the rendered page. We use slide tracking (<3->) 
    % to control exactly when individual arrows appear.
    \begin{tikzpicture}[remember picture, overlay, >=Stealth, thick, cyan]
        
        % Slide 3: Show the first alignment arrow (Subgroup -> Structure definition)
        \path<4-> [draw, ->]
            (ex1_txt1.north)
            [bend right=30] to
            node[below, sloped, text=black, font=\tiny]
                {Type being defined}
            (ex1_ln1.north);
        
        % Slide 4: Show the second alignment arrow with a descriptive label
        \path<5-> [draw, ->, draw=red]
            (ex1_txt2.west)
            [bend left=20] to
            node[below, sloped, text=black, font=\tiny]
                {Inductive Definition}
            (ex1_ln2.east);
        
    \end{tikzpicture}
\end{frame}

% Slide 9: Mechanization of Stmt
\begin{frame}[fragile]{The 6-Step Program to Profit}
    \begin{columns}[T]
        \begin{column}{0.55\textwidth}
            \begin{block}{Method}
\begin{enumerate}
    % Step 1: Target (Visible on slide 1, fades slightly or stays on later ones)
    \item Find something to mechanize
    % Step 2: Mechanize it (Visible on slide 2+)
    \item Then mechanize it
    % Step 3: Side-by-side comparisons (Visible on slide 3+)
    \item Put them side-by-side
    \item Repeat Steps 1-3
\end{enumerate}
            \end{block}
        \end{column}
        \begin{column}{0.45\textwidth}
            \begin{block}{Checklist}
\begin{itemize}
    \only<1-6>{\item Inductive types \No}
    \only<7>{\item Inductive types \checkmark}
    \item Semantic Functions \No
    \item Inference Rules \No
    \item Proofs and Lemmas \No
\end{itemize}
            \end{block}
        \end{column}
    \end{columns}

    % \vspace{0.5cm}
    
    % We wrap the side-by-side view inside an overprint or conditional block 
    % so it only shows up from Slide 1 onwards, changing state dynamically.
    \begin{columns}[T]
        
        % --- LEFT COLUMN: Lean Code ---
        \begin{column}{0.50\textwidth}
            \begin{block}{Lean 4 Formalization}
                \begin{lstlisting}[style=lean4, basicstyle=\footnotesize\ttfamily]
inductive Stm§\tikzmarknode{ex3_ln1}{t}§ where
  | ass (x : Var)§\tikzmarknode{ex3_ln2}{}§
        (a : Aexp)
  | skip §\tikzmarknode{ex3_ln3}{}§
  | composition (S§$_{1}$§ : Stmt)§\tikzmarknode{ex3_ln4}{}§
                (S§$_{2}$§ : Stmt)
  | §if§ (b : Bexp) §\tikzmarknode{ex3_ln5}{}§
       (S§$_{1}$§ S§$_{2}$§ : Stmt)
  | ...
                \end{lstlisting}
            \end{block}
        \end{column}

        % --- RIGHT COLUMN: Original Text Target---
        \begin{column}{0.48\textwidth}
            % Only show left column from slide 1, 3, and 4 (Skip slide 2 to isolate Lean)
            \only{
                \begin{block}{Original Definition}
\footnotesize
\[
\begin{array}{lcl}
    % \\[0.3em]
    \tikztarget{ex3_txt1}{S} &::=&\tikztarget{ex3_txt2}{ }x := a \\
    \\
    \\
      &\tikztarget{ex3_txt3}{ }\mid& \mathtt{skip} \\
      &\tikztarget{ex3_txt4}{ }\mid& S_1; S_2 \\
      \\
      &\tikztarget{ex3_txt5}{ }\mid& \mathtt{if}\ b\ \mathtt{then}\ S_1\ \mathtt{else}\ S_2 \\
      \\
      &\mid& \mathtt{while}\ b\ \mathtt{do}\ S
\end{array}
\]
                \end{block}
            }
        \end{column}
        
    \end{columns}

    % --- TIKZ OVERLAY LAYER ---
    % This executes "on top" of the rendered page. We use slide tracking (<3->) 
    % to control exactly when individual arrows appear.
    \begin{tikzpicture}[remember picture, overlay, >=Stealth, thick, cyan]
        
        \path<2> [draw, ->]
            (ex3_txt1.north)
            [bend right=30] to
            node[below, sloped, text=black, font=\tiny]
                {Type being defined}
            (ex3_ln1.north);
        
        \path<3> [draw, ->, draw=red]
            (ex3_txt2.west)
            [bend left=20] to
            node[above, sloped, text=black, font=\tiny]
                {Inductive Definition}
            (ex3_ln2.east);
            
        \path<4> [draw, ->, draw=red]
            (ex3_txt3.west)
            [bend right=5] to
            (ex3_ln3.east);
            
        \path<5> [draw, ->, draw=red]
            (ex3_txt4.west)
            [bend right=5] to
            (ex3_ln4.east);
            
        \path<6> [draw, ->, draw=red]
            (ex3_txt5.west)
            [bend right=5] to
            (ex3_ln5.east);
        
    \end{tikzpicture}
\end{frame}

% Slide 10: Transition from Types to Functions
\begin{frame}{So far so good}
    Perfect mapping from textbook BNF form \\
    into Lean 4 inductive typing \\ \vspace{1em}
    But what about the:
    \begin{itemize}
        \item \emph{Semantic Functions}
        \item \emph{Inference Rules}
        \item \emph{Proofs}
    \end{itemize}
    \vspace{1em}
    Well ...
        
\end{frame}

% Slide 11: Semantic functions Aexp_eval
\begin{frame}[fragile]{The 6-Step Program: \alert{Semantic Functions}}
    \begin{columns}[T]
        \begin{column}{0.55\textwidth}
            \begin{block}{Method}
\begin{enumerate}
    \item<1-> Find something to mechanize
    \item<2-> Then mechanize it
    \item<3-> Put them side-by-side
    \item<5-> Repeat Steps 1-3
\end{enumerate}
            \end{block}
        \end{column}
        \begin{column}{0.45\textwidth}
            \begin{block}{Checklist}
\begin{itemize}
    \item Inductive types \checkmark
    \only<1-4>{\item Semantic Functions \No}
    \only<5>{\item Semantic Functions \checkmark}
    \item Inference Rules \No
    \item Proofs and Lemmas \No
\end{itemize}
            \end{block}
        \end{column}
    \end{columns}

    \begin{columns}[T]
        
        % --- LEFT COLUMN: Lean Code ---
        \begin{column}{0.50\textwidth}
            \begin{block}<2->{Lean 4 Formalization}
                \begin{lstlisting}[style=lean4, basicstyle=\scriptsize\ttfamily]
def Aexp_eval :
  Aexp §\(\to\)§ (State §\(\to\)§ §\(\mathbb{Z}\)§)
  | .num §\tikzmarknode{ex4_ln1}{n}§, _ => §\(\mathcal{N}[\![\)§n§\(]\!]\)§
  | .var x, s => s x
  | §\(a_1 + a_2\)§, s => §\(\mathcal{A}[\![a_1]\!]\)§s + §\(\mathcal{A}[\![a_2]\!]\)§s
  | §\(a_1 \star a_2\)§, s => §\(\mathcal{A}[\![a_1]\!]\)§s * §\(\mathcal{A}[\![a_2]\!]\)§s
  | §\(a_1 - a_2\)§, s => §\(\mathcal{A}[\![a_1]\!]\)§s - §\(\mathcal{A}[\![a_2]\!]\)§s
                \end{lstlisting}
            \end{block}
        \end{column}

        % --- RIGHT COLUMN: Original Text Target---
        \begin{column}{0.48\textwidth}
            \begin{block}<1,3->{Original Definition}
\scriptsize
\[
\begin{array}{lcl}
        \mathcal{A}: Aexp \to \mathbb{Z}
        \\[0.8em]
        \mathcal{A}[\![\tikztarget{ex4_txt1}{n}]\!]s &=& \mathcal{N}[\![n]\!] \\
        \mathcal{A}[\![x]\!]s &=& s \, x \\
        \mathcal{A}[\![a_1 + a_2]\!]s &=& \mathcal{A}[\![a_1]\!]s + \mathcal{A}[\![a_2]\!]s \\
        \mathcal{A}[\![a_1 \star a_2]\!]s &=& \mathcal{A}[\![a_1]\!]s \cdot \mathcal{A}[\![a_2]\!]s \\
        \mathcal{A}[\![a_1 - a_2]\!]s &=& \mathcal{A}[\![a_1]\!]s - \mathcal{A}[\![a_2]\!]s
\end{array}
\]
            \end{block}
        \end{column}
        
    \end{columns}

    \begin{tikzpicture}[remember picture, overlay, >=Stealth, thick, cyan]
        
        \path<4> [draw, ->]
            (ex4_txt1.north)
            [bend right=30] to
            node[below, sloped, text=black, font=\tiny]
                {Type being defined}
            (ex4_ln1.north);
        
    \end{tikzpicture}


    \begin{block}<3>
      {What about inference rules?}
    \end{block}
\end{frame}

% Slide 12: Inference rules 1/3
\begin{frame}[fragile]{The 6-Step Program: \alert{Inference Rules} (NS)}
    \begin{columns}[T]
        \begin{column}{0.55\textwidth}
            \begin{block}{Method}
\begin{enumerate}
    \item<1-> Find something to mechanize
    \item<2-> Then mechanize it
    \item<3-> Put them side-by-side
    \item<5-> Repeat Steps 1-3
\end{enumerate}
            \end{block}
        \end{column}
        \begin{column}{0.45\textwidth}
            \begin{block}{Checklist}
\begin{itemize}
    \item Inductive types \checkmark
    \item Semantic Functions \checkmark
    \only<1-7>{\item Inference Rules \No}
    \only<8>{\item Inference Rules \checkmark}
    \item Proofs and Lemmas \No
\end{itemize}
            \end{block}
        \end{column}
    \end{columns}

    \begin{columns}[T]
        
        % --- LEFT COLUMN: Lean Code ---
        \begin{column}{0.50\textwidth}
            \begin{block}{Lean 4 Formalization}
            % \onslide<2->{
                \begin{lstlisting}[style=lean4, basicstyle=\scriptsize\ttfamily]
inductive big_step : Stmt §\(\to\)§
  State §\(\to\)§ State §\(\to\)§ Prop where
  | ass {s x a} : §\tikzmarknode{ex5_ln1}{}§
      §\(\langle\)§x §\(: \equiv\)§ a, s§\(\rangle\)§ §\(\to_{ns}\)§ §\tikzmarknode{ex5_ln2}{}§
      (s[x §\mapsto§ §\(\mathcal{A}[\![a]\!]\)§ s]) §\tikzmarknode{ex5_ln3}{}§
                \end{lstlisting}
            % }
            % \onslide<5->{
                \begin{lstlisting}[style=lean4, basicstyle=\scriptsize\ttfamily]
  | if_true {b S§$_{1}$§ S§$_{2}$§ s s'}
    (h_cond : §\(\mathcal{B}[\![b]\!]\)§ s = §\tikzmarknode{ex5_ln5}{tt}§)
    (h_then : §\(\langle\)§S§\(_1\)§, s§\(\rangle\)§ §\(\to_{ns}\)§ §\tikzmarknode{ex5_ln4}{s'}§) :
    §\(\langle\)§§if§ b §then§ S§$_{1}$§ §else§ S§$_{2}$§, s§\(\rangle\)§ §\(\to_{ns}\)§ s' §\tikzmarknode{ex5_ln6}{}§
                \end{lstlisting}
            % }
            \end{block}
        \end{column}

        % --- RIGHT COLUMN: Original Text Target---
        \begin{column}{0.48\textwidth}
            % Only show left column from slide 1, 3, and 4 (Skip slide 2 to isolate Lean)
            \begin{block}{Original Definition}
\vspace{2em}
\scriptsize
\[
  \begin{prooftree}
    \hypo{} % Empty hypothesis for an axiom
    \infer[
      left label=\tikzmarknode{ex5_txt1}{}[ass$_{\text{ns}}$]
      ]1{
          \parbox{2.5cm}{$
      \tikzmarknode{ex5_txt2}{}\langle x := a, s \rangle \to \\
      \tikzmarknode{ex5_txt3}{}s[x \mapsto \mathcal{A}[\![a]\!] s]
          $}
      }
  \end{prooftree}
\]
\vspace{2em}
\[
  \begin{prooftree}
    \hypo{\tikzmarknode{ex5_txt4}{}\langle S_1, s \rangle \to s'}
    \infer[
      left label={
        \parbox{0.8cm}{
          [if$_{\text{ns}}^{\text{\ tt}}$]\\
          \tikzmarknode{ex5_txt5}{}$\mathcal{B}[\![b]\!]s$ \\
          $= \mathbf{tt}$}
        }
      ]1{
        \tikzmarknode{ex5_txt6}{}\langle \mathtt{if}\ b\ \mathtt{then}\ S_1\ \mathtt{else}\ S_2, s \rangle \to s'
      }
  \end{prooftree}
\]
% \[
%   \begin{prooftree}
%     \hypo{\langle S_1, s \rangle \to s'}
%     \hypo{\langle S_2, s' \rangle \to s''}
%     \infer[
%       left label=[comp$_{\text{ns}}$]
%       ]2{
%           \langle S_1;\ S_2, s \rangle \to s''
%       }
%   \end{prooftree}
% \]
            \end{block}
        \end{column}
    \end{columns}


    \begin{tikzpicture}[remember picture, overlay, >=Stealth, thick, cyan]
        
        \path<2-4> [draw, ->]
            (ex5_txt1.west)
            to
            (ex5_ln1.east);
        \path<3-4> [draw, ->]
            (ex5_txt2.west)
            to
            (ex5_ln2.east);
        \path<4> [draw, ->]
            (ex5_txt3.west)
            to
            (ex5_ln3.east);

        \path<5-7> [draw, ->, draw=red]
            (ex5_txt4.west)
            [bend right=45] to
            (ex5_ln4.east);
        \path<6-7> [draw, ->, draw=red]
            (ex5_txt5.west)
            [bend right=15] to
            (ex5_ln5.east);
        \path<7-7> [draw, ->, draw=red]
            (ex5_txt6.west)
            [bend left=20] to
            (ex5_ln6.east);
        
    \end{tikzpicture}
\end{frame}

% Slide 13: Lemma 2.5 definition
\begin{frame}[fragile]{The 6-Step Program: \alert{Proofs} (NS) 1/10}
    \begin{columns}[T]
        \begin{column}{0.55\textwidth}
        \end{column}
        \begin{column}{0.45\textwidth}
            \begin{block}{Checklist}
\begin{itemize}
    \item Inductive types \checkmark
    \item Semantic Functions \checkmark
    \item Inference Rules \checkmark
    \item Proofs and Lemmas \No
\end{itemize}
            \end{block}
        \end{column}
    \end{columns}

    \begin{columns}[T]
        
        % --- LEFT COLUMN: Lean Code ---
        \begin{column}{0.50\textwidth}
            \begin{block}{Lean 4 Formalization}
                \begin{lstlisting}[style=lean4, basicstyle=\scriptsize\ttfamily]
theorem lemma_2_5
  (b : Bexp) (S : Stmt)
  (s s' : State) :
  (§\(\langle\)§§while§ b §then§ S, s§\(\rangle\)§ §\(\to_{ns}\)§ s')
  §\(\leftrightarrow\)§ 
  (§\(\langle\)§§if§ b
    §then§ (S; §while§ b §then§ S)
    §else§ Stmt.skip, s§\(\rangle\)§
  §\(\to_{ns}\)§ s') := by
                \end{lstlisting}
            \end{block}
        \end{column}

        % --- RIGHT COLUMN: Original Text Target---
        \begin{column}{0.48\textwidth}
            % Second half of definition
            \only{
                \begin{block}{Original Definition}
\vspace{0.4em}
\small
\textbf{Lemma 2.5}\\
\vspace{0.4em}
\scriptsize
The statement
\[\mathtt{while}\ b\ \mathtt{do}\ S\]
is semantically equivalent to
\[\mathtt{if}\ b\ \mathtt{then} (S;\ \mathtt{while}\ b\ \mathtt{do}\ S)\ \mathtt{else}\ \mathtt{skip}\]
                \end{block}
            }
        \end{column}
        
    \end{columns}
\end{frame}

% Slide 14: Lemma 2.5 equiv stmt
\begin{frame}[fragile]{The 6-Step Program: \alert{Proofs} (NS) 2/10}
    \begin{columns}[T]
        \begin{column}{0.55\textwidth}
        \end{column}
        \begin{column}{0.45\textwidth}
            \begin{block}{Checklist}
\begin{itemize}
    \item Inductive types \checkmark
    \item Semantic Functions \checkmark
    \item Inference Rules \checkmark
    \item Proofs and Lemmas \No
\end{itemize}
            \end{block}
        \end{column}
    \end{columns}

    \begin{columns}[T]
        
        % --- LEFT COLUMN: Lean Code ---
        \begin{column}{0.50\textwidth}
            \begin{block}{Lean 4 Formalization}
                \begin{lstlisting}[style=lean4, basicstyle=\scriptsize\ttfamily]
  apply Iff.intro

  -- Forward direction:
  -- if => while
  case mp =>
    exact if_compact b S s s'

  -- Backward direction:
  -- while => if
  case mpr =>
    exact while_expand b S s s'
                \end{lstlisting}
            \end{block}
        \end{column}

        % --- RIGHT COLUMN: Original Text Target---
        \begin{column}{0.48\textwidth}
            % Second half of definition
            \begin{block}{Original Definition}
\vspace{0.8em}
\scriptsize
The proof is in two parts
            \end{block}
        \end{column}
    \end{columns}
\end{frame}

% Slide 15: Lemma 2.5 while_expand 1/7
\begin{frame}[fragile]{The 6-Step Program: \alert{Proofs} (NS) 3/10}
    \begin{columns}[T]
        \begin{column}{0.55\textwidth}
        \end{column}
        \begin{column}{0.45\textwidth}
            \begin{block}{Checklist}
\begin{itemize}
    \item Inductive types \checkmark
    \item Semantic Functions \checkmark
    \item Inference Rules \checkmark
    \item Proofs and Lemmas \No
\end{itemize}
            \end{block}
        \end{column}
    \end{columns}
    
    \begin{columns}[T]
        
        % --- LEFT COLUMN: Lean Code ---
        \begin{column}{0.50\textwidth}
            \begin{block}{Lean 4 Formalization}
                \begin{lstlisting}[style=lean4, basicstyle=\scriptsize\ttfamily]
theorem while_expand
  (b : Bexp) (S : Stmt)
  (s s'' : State) :
  (§\(\langle\)§§while§ b §then§ S, s§\(\rangle\)§ §\(\to_{ns}\)§ s'') §\tikzmarknode{ex6_ln1}{}§
  §\(\to\)§ §\tikzmarknode{ex6_ln2}{}§
  (§\(\langle\)§§if§ b
    §then§ (S; §while§ b §then§ S) §\tikzmarknode{ex6_ln3}{}§
    §else§ Stmt.skip, s§\(\rangle\)§
  §\(\to_{ns}\)§ s'') := by
  intro h_while §\tikzmarknode{ex6_ln4}{}§
  cases h_while with §\tikzmarknode{ex6_ln5}{}§
                \end{lstlisting}
            \end{block}
        \end{column}

        % --- RIGHT COLUMN: Original Text Target---
        \begin{column}{0.48\textwidth}
            % Second half of definition
            \only{
                \begin{block}{Original Definition}
\vspace{0.8em}
\scriptsize
We shall first prove that if
\[
\tikzmarknode{ex6_txt1}{} \langle \mathtt{while}\ b\ \mathtt{do}\ S,\ s\rangle \to s'' \tag{*}
\]
\tikzmarknode{ex6_txt2}{}then
\[
  \parbox{4cm}{
    \tikzmarknode{ex6_txt3}{} $\langle \mathtt{if}\ b\ \mathtt{then} (S;\ \mathtt{while}\ b\ \mathtt{do}\ S)$\\
    $\mathtt{else}\ \mathtt{skip},\ s\rangle \to s''$
  } \tag{**}
\]
\tikzmarknode{ex6_txt4}{} Because (*) holds, we know that we have a derivation tree T for it.\\
\tikzmarknode{ex6_txt5}{}It can have one of two forms depending on whether it has been constructed using the
rule [while$_{\text{ns}}^{\text{\ tt}}$] or the axiom [while$_{\text{ns}}^{\text{\ ff}}$].
                \end{block}
            }
        \end{column}
        
    \end{columns}


    \begin{tikzpicture}[remember picture, overlay, >=Stealth, thick, cyan]
        
        \path<2> [draw, ->]
            (ex6_txt1.west)
            [bend right=5] to
            (ex6_ln1.east);

        \path<3> [draw, ->]
            (ex6_txt2.west)
            [bend right=5] to
            (ex6_ln2.east);

        \path<4> [draw, ->]
            (ex6_txt3.west)
            [bend right=5] to
            (ex6_ln3.east);

        \path<5> [draw, ->]
            (ex6_txt4.west)
            [bend right=5] to
            (ex6_ln4.east);

        \path<6> [draw, ->]
            (ex6_txt5.west)
            [bend left=10] to
            (ex6_ln5.east);
        
    \end{tikzpicture}
\end{frame}

% Slide 16: Lemma 2.5 while_expand 2/7
\begin{frame}[fragile]{The 6-Step Program: \alert{Proofs} (NS) 4/10}
    \begin{columns}[T]
        \begin{column}{0.55\textwidth}
        \end{column}
        \begin{column}{0.45\textwidth}
            \begin{block}{Checklist}
\begin{itemize}
    \item Inductive types \checkmark
    \item Semantic Functions \checkmark
    \item Inference Rules \checkmark
    \item Proofs and Lemmas \No
\end{itemize}
            \end{block}
        \end{column}
    \end{columns}
    
    \begin{columns}[T]
        
        % --- LEFT COLUMN: Lean Code ---
        \begin{column}{0.50\textwidth}
            \begin{block}{Lean 4 Formalization}
                \begin{lstlisting}[style=lean4, basicstyle=\scriptsize\ttfamily]
  | while_true
    h_cond_true T§$_1$§ T§$_2$§ =>
                \end{lstlisting}
            \end{block}
        \end{column}

        % --- RIGHT COLUMN: Original Text Target---
        \begin{column}{0.48\textwidth}
            % Second half of definition
            \only{
                \begin{block}{Original Definition}
\vspace{0.8em}
\scriptsize
In the first case, the derivation tree T has the form
\[
    \begin{prooftree}
        \hypo{T_1}
        \hypo{T_2}
        \infer[
            ]2{
            \langle \mathtt{while}\ b\ \mathtt{do}\ S,\ s\rangle \to s''
            }
    \end{prooftree}
\]
where \(T_1\) is a derivation tree with root \(\langle S,\ s\rangle \to s'\) and \(T_2\) is a derivation tree
with root \(\langle \mathtt{while}\ b\ \mathtt{do}\ S,\ s'\rangle \to s''\).
Furthermore, \(\mathcal{B}[\![b]\!]s = \mathbf{tt}\).
                \end{block}
            }
        \end{column}
        
    \end{columns}
\end{frame}

% Slide 17: Lemma 2.5 while_expand 3/7
\begin{frame}[fragile]{The 6-Step Program: \alert{Proofs} (NS) 5/10}
    \begin{columns}[T]
        \begin{column}{0.55\textwidth}
        \end{column}
        \begin{column}{0.45\textwidth}
            \begin{block}{Checklist}
\begin{itemize}
    \item Inductive types \checkmark
    \item Semantic Functions \checkmark
    \item Inference Rules \checkmark
    \item Proofs and Lemmas \No
\end{itemize}
            \end{block}
        \end{column}
    \end{columns}
    
    \begin{columns}[T]
        
        % --- LEFT COLUMN: Lean Code ---
        \begin{column}{0.50\textwidth}
            \begin{block}{Lean 4 Formalization}
                \begin{lstlisting}[style=lean4, basicstyle=\scriptsize\ttfamily]
    have h_comp :
      (§\(\langle\)§S; while b then S, s§\(\rangle\)§
      §\(\to_{ns}\)§ s'') :=
        big_step.comp T§$_1$§ T§$_2$§
                \end{lstlisting}
            \end{block}
        \end{column}

        % --- RIGHT COLUMN: Original Text Target---
        \begin{column}{0.48\textwidth}
            % Second half of definition
            \only{
                \begin{block}{Original Definition}
\vspace{0.8em}
\scriptsize
Using the derivation trees \(T_1\) and \(T_2\) as the premises for the rules [comp$_{\text{ns}}$], we can construct the derivation tree
\[
    \begin{prooftree}
        \hypo{T_1}
        \hypo{T_2}
        \infer[
            % left label=[comp$_{\text{ns}}$]
            ]2{
            \langle S;\ \mathtt{while}\ b\ \mathtt{do}\ S,\ s\rangle \to s''
            }
    \end{prooftree}
\]
                \end{block}
            }
        \end{column}
        
    \end{columns}
\end{frame}

% Slide 18: Lemma 2.5 while_expand 4/7
\begin{frame}[fragile]{The 6-Step Program: \alert{Proofs} (NS) 6/10}
    \begin{columns}[T]
        \begin{column}{0.55\textwidth}
        \end{column}
        \begin{column}{0.45\textwidth}
            \begin{block}{Checklist}
\begin{itemize}
    \item Inductive types \checkmark
    \item Semantic Functions \checkmark
    \item Inference Rules \checkmark
    \item Proofs and Lemmas \No
\end{itemize}
            \end{block}
        \end{column}
    \end{columns}
    
    \begin{columns}[T]
        
        % --- LEFT COLUMN: Lean Code ---
        \begin{column}{0.50\textwidth}
            \begin{block}{Lean 4 Formalization}
                \begin{lstlisting}[style=lean4, basicstyle=\scriptsize\ttfamily]
    exact big_step.if_true
      h_cond_true -- Condition
      h_comp      -- Then case
                \end{lstlisting}
            \end{block}
        \end{column}

        % --- RIGHT COLUMN: Original Text Target---
        \begin{column}{0.48\textwidth}
            % Second half of definition
            \only{
                \begin{block}{Original Definition}
\vspace{0.8em}
\scriptsize
Using that \(\mathcal{B}[\![b]\!]s = \mathbf{tt}\), we can use the rule [if$_{\text{ns}}^{\ \text{tt}}$]
to construct the derivation tree
\[
  \begin{prooftree}
    \hypo{T_1}
    \hypo{T_2}
    \infer2{
      \langle S;\ \mathtt{while}\ b\ \mathtt{do}\ S,\ s\rangle \to s''
    }
    \infer1{
      \parbox{4cm}{$
        \langle \mathtt{if}\ b\ \mathtt{then}\ (S;\ \mathtt{while}\ b\ \mathtt{do}\ S)\ \\
        \mathtt{else}\ \mathtt{skip},\ s\rangle \to s''
      $}
    }
  \end{prooftree}
\]
hereby showing that (**) holds.
                \end{block}
            }
        \end{column}
        
    \end{columns}
\end{frame}

% Slide 19: Lemma 2.5 while_expand 5/7
\begin{frame}[fragile]{The 6-Step Program: \alert{Proofs} (NS) 7/10}
    \begin{columns}[T]
        \begin{column}{0.55\textwidth}
        \end{column}
        \begin{column}{0.45\textwidth}
            \begin{block}{Checklist}
\begin{itemize}
    \item Inductive types \checkmark
    \item Semantic Functions \checkmark
    \item Inference Rules \checkmark
    \item Proofs and Lemmas \No
\end{itemize}
            \end{block}
        \end{column}
    \end{columns}
    
    \begin{columns}[T]
        
        % --- LEFT COLUMN: Lean Code ---
        \begin{column}{0.50\textwidth}
            \begin{block}{Lean 4 Formalization}
                \begin{lstlisting}[style=lean4, basicstyle=\scriptsize\ttfamily]
  | while_false h_cond_false =>
                \end{lstlisting}
            \end{block}
        \end{column}

        % --- RIGHT COLUMN: Original Text Target---
        \begin{column}{0.48\textwidth}
            % Second half of definition
            \only{
                \begin{block}{Original Definition}
\vspace{0.8em}
\scriptsize
Alternatively, the derivation tree T is an instance of [while$_{\text{ns}}^{\ \text{ff}}$].
Then \(\mathcal{B}[\![b]\!]s = \mathbf{ff}\) and we must have that \(s'' = s\). So T simply is
\[
    \begin{prooftree}
        \hypo{\langle \mathtt{while}\ b\ \mathtt{do}\ S,\ s\rangle \to s}
    \end{prooftree}
\]
                \end{block}
            }
        \end{column}
        
    \end{columns}
\end{frame}

% Slide 20: Lemma 2.5 while_expand 6/7
\begin{frame}[fragile]{The 6-Step Program: \alert{Proofs} (NS) 8/10}
    \begin{columns}[T]
        \begin{column}{0.55\textwidth}
        \end{column}
        \begin{column}{0.45\textwidth}
            \begin{block}{Checklist}
\begin{itemize}
    \item Inductive types \checkmark
    \item Semantic Functions \checkmark
    \item Inference Rules \checkmark
    \item Proofs and Lemmas \No
\end{itemize}
            \end{block}
        \end{column}
    \end{columns}
    
    \begin{columns}[T]
        
        % --- LEFT COLUMN: Lean Code ---
        \begin{column}{0.50\textwidth}
            \begin{block}{Lean 4 Formalization}
                \begin{lstlisting}[style=lean4, basicstyle=\scriptsize\ttfamily]
    have skip_stmt :
      (§\(\langle\)§Stmt.skip, s§\(\rangle\)§
      §\(\to_{ns}\)§ s) :=
      big_step.skip
                \end{lstlisting}
            \end{block}
        \end{column}

        % --- RIGHT COLUMN: Original Text Target---
        \begin{column}{0.48\textwidth}
            % Second half of definition
            \only{
                \begin{block}{Original Definition}
\vspace{0.8em}
\scriptsize
Using the axiom [skip$_{\text{ns}}$], we get a derivation tree
\[
    \begin{prooftree}
        \hypo{\langle \mathtt{skip},\ s\rangle \to s''}
    \end{prooftree}
\]
                \end{block}
            }
        \end{column}
        
    \end{columns}
\end{frame}

% Slide 21: Lemma 2.5 while_expand 7/7
\begin{frame}[fragile]{The 6-Step Program: \alert{Proofs} (NS) 9/10}
    \begin{columns}[T]
        \begin{column}{0.55\textwidth}
        \end{column}
        \begin{column}{0.45\textwidth}
            \begin{block}{Checklist}
\begin{itemize}
    \item Inductive types \checkmark
    \item Semantic Functions \checkmark
    \item Inference Rules \checkmark
    \item Proofs and Lemmas \No
\end{itemize}
            \end{block}
        \end{column}
    \end{columns}
    
    \begin{columns}[T]
        
        % --- LEFT COLUMN: Lean Code ---
        \begin{column}{0.50\textwidth}
            \begin{block}{Lean 4 Formalization}
                \begin{lstlisting}[style=lean4, basicstyle=\scriptsize\ttfamily]
    exact big_step.if_false
      h_cond_false
      skip_stmt
                \end{lstlisting}
            \end{block}
        \end{column}

        % --- RIGHT COLUMN: Original Text Target---
        \begin{column}{0.48\textwidth}
            % Second half of definition
            \only{
                \begin{block}{Original Definition}
\vspace{0.8em}
\scriptsize
and we can now apply the rule [if$_{\text{ns}}^{\ \text{ff}}$] to construct a derivation tree for (**):
\[
  \begin{prooftree}
    \hypo{\langle \mathtt{skip},\ s\rangle \to s''}
    \infer1{
      \parbox{4cm}{$
        \langle \mathtt{if}\ b\ \mathtt{then}\ (S;\ \mathtt{while}\ b\ \mathtt{do}\ S)\ \\
        \mathtt{else}\ \mathtt{skip},\ s\rangle \to s''
      $}
    }
  \end{prooftree}
\]
This completes the first part of the proof.
                \end{block}
            }
        \end{column}
        
    \end{columns}
\end{frame}

% Slide 22: Lemma 2.5 final stmt
\begin{frame}[fragile]{The 6-Step Program: \alert{Proofs} (NS) 10/10}
    \begin{columns}[T]
        \begin{column}{0.55\textwidth}
        \end{column}
        \begin{column}{0.45\textwidth}
            \begin{block}{Checklist}
\begin{itemize}
    \item Inductive types \checkmark
    \item Semantic Functions \checkmark
    \item Inference Rules \checkmark
    \only<1>{\item Proofs and Lemmas \No}
    \only<2>{\item Proofs and Lemmas \checkmark}
\end{itemize}
            \end{block}
        \end{column}
    \end{columns}

    (**) to (*) (reverse direction)\\
    Same pattern:
    \begin{enumerate}
      \item Split into cases
      \item Use premises to construct new trees
      \item Tell lean that new tree is wanted tree
    \end{enumerate}
\end{frame}

% Conclusion time
\end{document}